\chapter{The IS--LM Model}\label{ch:b1:is-lm}

On 16 December 2008 the Federal Reserve set a target range of 0 to 1/4 per
cent for its policy rate, and it kept it there until 16 December 2015.
% ledger: b1:fed-zlb-2008
% ledger: b1:fed-liftoff-2015
Over those seven years the effective rate never rose above $0.22\%$, while
the Federal Reserve's balance sheet grew from 0.891 trillion dollars at the
end of 2007 to 4.498 trillion at the end of 2014. % figdata: zlb_stats
% figdata: assets_points
A central-bank balance sheet five times as large did not push the interest rate
below zero. This chapter builds the short-run model in which that is the
expected result: output set by demand, at a price level that has not yet
moved.

\begin{openingchart}{The effective federal funds rate (per cent, left axis,
blue) and the Federal Reserve's total assets (trillions of dollars, right
axis, red), monthly, 2003--2025.}{Board of Governors of the Federal
Reserve System, H.15 and H.4.1}
\begin{tikzpicture}
\begin{axis}[omtime, xlabel={Year}, ylabel={Federal funds rate (\%)}, xmin=2003, xmax=2026, ymin=0, ymax=6,
  width=0.72\linewidth, height=0.4\linewidth, axis y line*=left]
  \addplot+[thick, no marks, omDef] table[x=t, y=ffr, col sep=comma] {figdata/bachelor-1/24-is-lm/out/fedzlb.csv};
\end{axis}
\begin{axis}[omtime, xmin=2003, xmax=2026, ymin=0, ymax=10, width=0.72\linewidth, height=0.4\linewidth,
  axis y line*=right, axis x line=none, ylabel={Total assets (trillion \$)}, grid=none]
  \addplot+[thick, no marks, omThm] table[x=t, y=assets, col sep=comma] {figdata/bachelor-1/24-is-lm/out/fedzlb.csv};
\end{axis}
\end{tikzpicture}
\end{openingchart}

The rate fell to its floor in late 2008 and stayed there for seven years; it
fell there again in 2020. The balance sheet grew in steps during the first
episode, and to 8.939 trillion dollars in April 2022 after the second.
% figdata: assets_points
In Japan the short-term rate was $0.059\%$ in 1999, $0.473\%$ in 2007 and $0.073\%$
in 2015. % figdata: jst_rates

\begin{recall}
The Keynesian cross and the multiplier, with its limits, and the policy
rate were met at school. Money demand is
\cref{def:b1:money-and-inflation:money-demand}; the \omterm{def:b1:long-run-model:loanable-funds}{loanable funds market}
of the long run is \cref{def:b1:long-run-model:loanable-funds}; implicit
differentiation is \cref{prop:b1:comparative-statics:implicit}.
\end{recall}

\section{The goods market and the IS curve}\label{sec:b1:is-lm:is}

In the short run the price level $P$ is fixed and output is whatever is
demanded. Planned spending in a closed economy is
\[
E = C(Y - T) + I(r) + G,
\]
with $0 < C' < 1$ and $I' < 0$. The goods market is in \omterm{def:b1:models-and-data:equilibrium}{equilibrium} when output
equals planned spending, $Y = E$.

\begin{definition}[IS curve]\label{def:b1:is-lm:is}
\index{IS curve}
The \emph{IS curve} is the set of pairs $(Y, r)$ at which the goods market is
in \omterm{def:b1:models-and-data:equilibrium}{equilibrium}: $Y = C(Y - T) + I(r) + G$.
\end{definition}

Differentiating, $dY = C'dY + I'dr$, so along the curve
\[
\frac{dr}{dY} = \frac{1 - C'}{I'} < 0 :
\]
a lower interest rate raises investment, and output rises by the
multiplier $1/(1 - C')$ times the extra investment. A rise in $G$ shifts the
curve to the right by $\Delta G/(1 - C')$; a rise in $T$ shifts it to the left by
$C'\Delta T/(1 - C')$.

\section{The money market and the LM curve}\label{sec:b1:is-lm:lm}

With $P$ fixed and expected inflation zero, the nominal and real interest
rates coincide, and the money market clears when $M/P = L(r, Y)$, with
$L_r < 0 < L_Y$.

\begin{definition}[LM curve]\label{def:b1:is-lm:lm}
\index{LM curve}
The \emph{LM curve} is the set of pairs $(Y, r)$ at which the money market is
in \omterm{def:b1:models-and-data:equilibrium}{equilibrium} for a given real money supply: $M/P = L(r, Y)$.
\end{definition}

Along it, $0 = L_r dr + L_Y dY$, so $dr/dY = -L_Y/L_r > 0$: more output raises money
demand, and with a given money supply the interest rate must rise to
choke it off. A larger $M$ shifts the curve down.

\begin{measure}[A policy rate and a market rate]
The rate a central bank sets is an overnight rate at which banks lend
reserves to one another. The rates that matter for investment --- on loans,
bonds, mortgages --- add a premium for time and for risk. The model has a
single $r$; in the data the spread between the two moves, and in 2008 it
widened sharply.
\end{measure}

\section{Equilibrium and policy}\label{sec:b1:is-lm:equilibrium}

\begin{model}{The IS--LM model}
\index{IS--LM model}
\assumes A fixed price level; a closed economy; no expectations; a money
supply (LM) or an interest rate (MP) set by the central bank.
\predicts A fiscal expansion raises output and the interest rate (with
LM), or output alone (with MP); a monetary expansion lowers the interest
rate and raises output, except in a \omterm{def:b1:is-lm:liquidity-trap}{liquidity trap}.
\tested The signs of the responses to monetary shocks: identified from the
Federal Reserve's own records for 1969--1996, policy shocks have large,
fairly rapid effects on output and inflation (Romer and Romer 2004).
% ledger: b1:romer2004
At the \omterm{def:b1:is-lm:liquidity-trap}{zero lower bound}, more central-bank money left the policy rate at
its floor (opening chart).
\breaks Prices move (\cref{ch:b1:aggregate-supply}); economies are open
(\cref{ch:b1:open-economy}); expectations matter and finance can break, as
in 2008 --- subjects of the second university year. It is a toy model: its
signs are informative, the size of its multipliers is not to be trusted
(\cref{ch:b1:stabilisation-policy}).
\end{model}

\begin{proposition}[Policy multipliers]\label{prop:b1:is-lm:multipliers}
In the IS--LM model,
\[
\frac{dY}{dG} = \frac{1}{1 - C' + I'L_Y/L_r}, \qquad
\frac{dY}{d(M/P)} = \frac{I'/L_r}{1 - C' + I'L_Y/L_r}.
\]
Both are positive, and $0 < dY/dG < 1/(1 - C')$: \omterm{def:b1:long-run-model:crowding-out}{crowding out} is partial.
\end{proposition}

\begin{proof}
Differentiate the two equations: $(1 - C')dY - I'dr = dG$ and $L_Y dY + L_r dr = d(M/P)$.
From the second, $dr = (d(M/P) - L_Y dY)/L_r$; substituting,
$(1 - C' + I'L_Y/L_r)dY = dG + (I'/L_r)d(M/P)$. Since $I' < 0$ and $L_r < 0$, $I'L_Y/L_r > 0$, so the
denominator exceeds $1 - C' > 0$, and $I'/L_r > 0$.
\end{proof}

\begin{definition}[Fiscal multiplier]\label{def:b1:is-lm:fiscal-multiplier}
\index{fiscal multiplier}
The \emph{fiscal multiplier} is the change in output per unit change in
government purchases, $dY/dG$, all else given; in IS--LM it is smaller than
the Keynesian-cross multiplier $1/(1 - C')$ because the interest rate rises and
investment falls.
\end{definition}

\begin{omfigure}
\begin{tikzpicture}
\begin{axis}[omrecord, grid=none, xlabel={Output $Y$ (billion dalens)}, ylabel={Interest rate $r$ (\%)},
  xmin=340, xmax=420, ymin=-1, ymax=6, width=0.58\linewidth, height=0.42\linewidth]
  \addplot[demand, domain=340:420] {(430.765 - x)/12.5};
  \addplot[shifted, domain=340:420] {(405.765 - x)/12.5};
  \addplot[supply, domain=340:420] {0.1*x - 36.326};
  \addplot[guide] coordinates {(393.264,-1) (393.264,3)};
  \addplot[guide] coordinates {(382.153,-1) (382.153,1.889)};
  \node[font=\footnotesize, anchor=west] at (axis cs:407,1.25) {IS};
  \node[font=\footnotesize, anchor=north east] at (axis cs:356,3.5) {IS$'$};
  \node[font=\footnotesize, anchor=north west] at (axis cs:405,4.4) {LM};
\end{axis}
\end{tikzpicture}

{\small The weekend problem's economy: IS (blue) and LM (red) cross at
$Y = 393.264$ and $r = 3\%$. A fall of 10 billion in exports shifts IS to the left by
25 (dashed); the new \omterm{def:b1:models-and-data:equilibrium}{equilibrium} is $Y = 382.153$ and $r = 1.889\%$: the fall in the
rate saves part of the output.}
\end{omfigure}

With the weekend problem's numbers ($C' = 0.6$, $I' = -5$ billion per point,
$L_Y = 0.5$, $L_r = -5$) the Keynesian-cross multiplier is $2.5$ and the IS--LM
multiplier $1/(0.4 + 0.5) = 1.111$. % computed: code/bachelor-1/24-is-lm/exercises/solutions.py

\section{The liquidity trap}\label{sec:b1:is-lm:trap}

The nominal interest rate cannot fall much below zero: people can hold
currency, which pays zero. When the \omterm{def:b1:is-lm:is}{IS curve} crosses the \omterm{def:b1:is-lm:lm}{LM curve} at a
negative rate, the rate stops at the floor.

\begin{definition}[Liquidity trap, zero lower bound]\label{def:b1:is-lm:liquidity-trap}
\index{liquidity trap}\index{zero lower bound}
The \emph{zero lower bound} is the floor, at or near zero, below which the
nominal interest rate cannot be pushed. An economy is in a \emph{liquidity
trap} when the interest rate is at that floor: money and bonds are then
\omterm{def:b1:preferences-and-utility:perfect}{perfect substitutes}, and more money is held without lowering the rate.
\end{definition}

In a trap the \omterm{def:b1:is-lm:lm}{LM curve} is flat at the floor. A larger money supply moves
the flat part to the right and leaves the \omterm{def:b1:models-and-data:equilibrium}{equilibrium} unchanged; a fiscal
expansion raises output by the full Keynesian-cross multiplier, since the
rate does not rise and nothing is crowded out.

\begin{omfigure}
\begin{tikzpicture}
\begin{axis}[omrecord, grid=none, xlabel={Output $Y$ (billion dalens)}, ylabel={Interest rate $r$ (\%)},
  xmin=330, xmax=420, ymin=-1, ymax=5, width=0.58\linewidth, height=0.42\linewidth]
  \addplot[supply, domain=330:420] {max(0, 0.1*x - 36.326)};
  \addplot[demand, domain=330:420] {(430.765 - 25 - 50 - x)/12.5};
  \addplot[shifted, domain=330:420] {(430.765 - 25 - 50 + 7.5 - x)/12.5};
  \addplot[guide] coordinates {(355.764,-1) (355.764,0)};
  \node[font=\footnotesize, anchor=north] at (axis cs:355.764,-0.15) {$355.8$};
  \node[font=\footnotesize, anchor=south east] at (axis cs:418,4) {LM};
  \node[font=\footnotesize, anchor=south west] at (axis cs:331,1.5) {IS};
\end{axis}
\end{tikzpicture}

{\small A \omterm{def:b1:is-lm:liquidity-trap}{liquidity trap} in the weekend problem: after a fall in exports and
in investment demand, IS (blue) crosses the \omterm{def:b1:is-lm:lm}{LM curve} (red) on its flat part
at $r = 0$, at $Y = 355.764$. A rise of 3 billion in $G$ (dashed) raises output by
7.5, the full multiplier of $2.5$; more money would not move the flat part
of LM.}
\end{omfigure}

\begin{record}[The zero lower bound, 2008--2015]
On 16 December 2008 the Federal Reserve set the target range for its
policy rate at 0 to 1/4 per cent, and raised it to 1/4 to 1/2 per cent only
on 16 December 2015. % ledger: b1:fed-zlb-2008
% ledger: b1:fed-liftoff-2015
In between, the effective rate stayed between $0.07\%$ and $0.22\%$ for 84
months, while total assets, already 2.5 times their end-2007 level, doubled
again. % figdata: zlb_stats, assets_points
In Japan the short-term rate was below $0.5\%$ from 1999 on. % figdata: jst_rates
As the model predicts, more money did not lower the rate. Whether it
raised output through other channels --- longer-term rates, asset prices,
expectations --- is a question for the second university year.
\end{record}

\section{The version with an interest-rate rule}\label{sec:b1:is-lm:mp}

Central banks do not set the money supply; they set an interest rate and
supply whatever money is demanded at it.

\begin{definition}[MP curve, interest-rate rule]\label{def:b1:is-lm:mp}
\index{MP curve}\index{interest-rate rule}
An \emph{interest-rate rule} gives the interest rate the central bank sets
as a function of the state of the economy. The \emph{MP curve} is that rule in
the $(Y, r)$ plane; in its simplest form it is horizontal at the rate the
bank has chosen, $r = \bar r$.
\end{definition}

With a horizontal \omterm{def:b1:is-lm:mp}{MP curve}, a fiscal expansion raises output by the full
multiplier unless the bank raises its rate; a fall in investment demand
lowers output unless the bank cuts. The model then reads as central banks
describe themselves: shocks shift IS, the bank moves MP, and output
responds. It is the starting point of the models of the second university
year.

\begin{omfigure}
\begin{tikzpicture}
\begin{axis}[omrecord, grid=none, xlabel={Output $Y$}, ylabel={Interest rate $r$ (\%)},
  xmin=800, xmax=1300, ymin=0, ymax=8, width=0.56\linewidth, height=0.4\linewidth]
  \addplot[demand, domain=800:1300] {(1500 - x)/100};
  \addplot[shifted, domain=800:1300] {(1340 - x)/100};
  \addplot[supply, domain=800:1300] {5};
  \addplot[supply, dashed, domain=800:1300] {3.4};
  \addplot[guide] coordinates {(1000,0) (1000,5)};
  \node[font=\footnotesize, anchor=south east] at (axis cs:1290,5) {MP};
  \node[font=\footnotesize, anchor=south east] at (axis cs:1290,3.4) {MP$'$};
  \node[font=\footnotesize, anchor=south west] at (axis cs:1160,3.45) {IS};
  \node[font=\footnotesize, anchor=north east] at (axis cs:1150,1.85) {IS$'$};
\end{axis}
\end{tikzpicture}

{\small IS--MP with the exercises' numbers: a fall of 40 in investment demand
shifts IS from $Y = 1500 - 100r$ to $Y = 1340 - 100r$ (dashed). Keeping output at 1,000
requires a rate cut from $5\%$ to $3.4\%$.}
\end{omfigure}

\section{The toy model in its frame}\label{sec:b1:is-lm:frame}

IS--LM holds the price level fixed, ignores what people expect, and treats
the economy as closed. Its predictions are signs: a fiscal expansion raises
output; a monetary expansion raises output except in a trap; a fall in
private demand lowers output unless policy offsets it. The signs survive
much better than the sizes. Its multipliers depend on slopes that are hard
to measure and that change when expectations do; the evidence on \omterm{def:b1:is-lm:fiscal-multiplier}{fiscal
multipliers} is the subject of the stabilisation chapter.

\section*{Exercises}\label{sec:b1:is-lm:exercises}

\begin{exercise}[$\star$]\label{exo:b1:is-lm:is-ex}
With $C = 50 + 0.75(Y - T)$, $I = 300 - 25r$ and $G = T = 100$, derive the \omterm{def:b1:is-lm:is}{IS curve}.
\end{exercise}

\begin{exercise}[$\star$]\label{exo:b1:is-lm:lm-ex}
With $M/P = 500$ and $L = Y - 100r$, derive the \omterm{def:b1:is-lm:lm}{LM curve}.
\end{exercise}

\begin{exercise}[$\star$]\label{exo:b1:is-lm:read}
The short-term rate was $5.02\%$ in the United States and $3.859\%$ in Germany in
2007, and $0.13\%$ and $-0.108\%$ in 2015. % figdata: jst_rates
Which were at or near the \omterm{def:b1:is-lm:liquidity-trap}{zero lower bound}? What does a negative rate say
about the bound?
\end{exercise}

\begin{exercise}[$\star$]\label{exo:b1:is-lm:sign}
Predict the sign of the change in $r$ after a rise in $G$ with an \omterm{def:b1:is-lm:lm}{LM curve} and
with an \omterm{def:b1:is-lm:mp}{MP curve}.
\end{exercise}

\begin{exercise}[$\star\star$]\label{exo:b1:is-lm:eq}
Find the \omterm{def:b1:models-and-data:equilibrium}{equilibrium} of \cref{exo:b1:is-lm:is-ex,exo:b1:is-lm:lm-ex} and the \omterm{def:b1:is-lm:fiscal-multiplier}{fiscal
multiplier}.
\end{exercise}

\begin{exercise}[$\star\star$]\label{exo:b1:is-lm:monetary}
In the same economy, find the monetary multiplier $dY/d(M/P)$.
\end{exercise}

\begin{exercise}[$\star\star$]\label{exo:b1:is-lm:trap-ex}
Investment demand collapses and the \omterm{def:b1:is-lm:is}{IS curve} becomes $Y = 400 - 100r$; the
interest rate cannot fall below zero. Find the \omterm{def:b1:models-and-data:equilibrium}{equilibrium}, and the effect
of doubling $M$.
\end{exercise}

\begin{exercise}[$\star\star$]\label{exo:b1:is-lm:mp-ex}
In the economy of \cref{exo:b1:is-lm:is-ex}, the central bank sets $r = 5$. A fall
in confidence lowers investment demand by 40 at every rate. Find the rate
cut that keeps output at 1,000.
\end{exercise}

\begin{exercise}[$\star\star$]\label{exo:b1:is-lm:assets}
The Federal Reserve's total assets were 0.891 trillion dollars at the end of
2007, 2.239 trillion at the end of 2008 and 4.487 trillion at the end of
2015. % figdata: assets_points
Compute the growth factors, and say what the LM model predicts for the
interest rate.
\end{exercise}

\begin{exercise}[$\star\star\star$]\label{exo:b1:is-lm:general}
Derive \cref{prop:b1:is-lm:multipliers} by implicit differentiation, and the
effect of $G$ on $r$.
\end{exercise}

\begin{exercise}[$\star\star\star$]\label{exo:b1:is-lm:crowding}
Using the multipliers, show when \omterm{def:b1:long-run-model:crowding-out}{crowding out} vanishes and when it is
complete, in terms of $L_r$ and $I'$.
\end{exercise}

\begin{exercise}[$\star\star\star$]\label{exo:b1:is-lm:refute}
What observation in 2008--2015 would have refuted the liquidity-trap
prediction?
\end{exercise}

\section{Problem: Brenland's Recession}\label{sec:b1:is-lm:problem}

\begin{problem}[{Weekend problem --- a recession in a toy model}]\label{pb:b1:is-lm:brenland}
Take Brenland's 2025 accounts (invented), in billions of dalens: $Y = 393.264$,
$C = 225.014$, $I = 88.484$, $G = 78.653$ and net exports $NX = 1.113$; the interest rate is
$r = 3\%$. Model them with $C = c_0 + 0.6(Y - T)$, $T = 80$ (a round figure for
the model, not the net taxes of the accounts, 70.788), $I = I_0 - 5r$, and money demand
$L = 0.5Y - 5r$; net exports are given. The price level is fixed and $r \ge 0$.

\medskip\noindent\textbf{Part I --- IS and LM.}
\begin{enumerate}
  \item Calibrate $c_0$ so that the consumption function fits 2025.
  \item Calibrate $I_0$.
  \item Find the real money supply that makes $r = 3$ an \omterm{def:b1:models-and-data:equilibrium}{equilibrium}.
  \item Write the \omterm{def:b1:is-lm:is}{IS curve}.
  \item Write the \omterm{def:b1:is-lm:lm}{LM curve}.
  \item Give the slopes $dr/dY$ of both curves.
  \item Compute the Keynesian-cross multiplier.
  \item Compute the IS--LM fiscal and monetary multipliers.
\end{enumerate}

\medskip\noindent\textbf{Part II --- A fall in exports.} Net exports fall by 10.
\begin{enumerate}[resume]
  \item By how much does the \omterm{def:b1:is-lm:is}{IS curve} shift?
  \item Find the new $Y$ and $r$.
  \item By how much does investment change, and why?
  \item What would the fall in output be with the interest rate held at $3\%$?
\end{enumerate}

\medskip\noindent\textbf{Part III --- Monetary and fiscal answers.}
\begin{enumerate}[resume]
  \item Find the rise in $M/P$ that restores output to 393.264.
  \item Find the interest rate then, and investment.
  \item Find instead the rise in $G$ that restores output.
  \item Find the interest rate then, and investment.
  \item Compare the two answers: what does each leave different?
\end{enumerate}

\medskip\noindent\textbf{Part IV --- The rate hits zero.} On top of the fall in
exports, investment demand falls by 20 at every rate.
\begin{enumerate}[resume]
  \item Find the \omterm{def:b1:models-and-data:equilibrium}{equilibrium} ignoring the floor.
  \item Find the \omterm{def:b1:models-and-data:equilibrium}{equilibrium} with the floor.
  \item What does a rise in $M/P$ do?
  \item Find the multiplier of a small rise in $G$.
  \item By how much can $G$ rise before the rate leaves zero?
  \item Find the rise in $G$ that restores output to 393.264.
  \item Why is the average multiplier of that rise smaller than 2.5?
  \item Which features of the toy model make the multiplier at the floor so large?
  \item State the result: what happens to the \omterm{def:b1:is-lm:fiscal-multiplier}{fiscal multiplier} at the zero bound, and why is that number the model's, not the economy's?
\end{enumerate}
\end{problem}
